\documentclass[11pt,a4paper]{article}
\usepackage[italian]{babel}
\usepackage[margin=2.5cm]{geometry}
\usepackage{parskip}
\usepackage{amsmath, amssymb, amsthm}
\usepackage{booktabs}
\usepackage{array}
\usepackage{xcolor}
\usepackage{needspace}
\usepackage{enumitem}
\usepackage{listings}
\usepackage{tikz}
\usepackage[most]{tcolorbox}
\usepackage[hidelinks]{hyperref}
\usetikzlibrary{decorations.pathreplacing, shapes.geometric, arrows.meta, positioning}

% Niente vedove/orfane: i paragrafi non lasciano righe isolate a fine/inizio pagina
\widowpenalty=10000
\clubpenalty=10000

% Elenchi più compatti
\setlist{itemsep=2pt, parsep=0pt, topsep=4pt, partopsep=0pt}

% --- Riquadri con bordo nero, senza numero (si spezzano tra due pagine) ---
% Uso: \begin{definizione} ... \end{definizione}
%      \begin{teorema}[Nome] ... \end{teorema}  (il nome è facoltativo)
\tcbset{nota/.style={enhanced jigsaw, breakable, colback=white,
  colframe=black, boxrule=0.5pt, sharp corners}}
\NewTColorBox{definizione}{}{nota,
  before upper={\textbf{Definizione.}\ }}
\NewTColorBox{teorema}{o}{nota,
  before upper={\textbf{Teorema\IfValueT{#1}{ (#1)}.}\ }}
\NewTColorBox{proposizione}{o}{nota,
  before upper={\textbf{Proposizione\IfValueT{#1}{ (#1)}.}\ }}
\NewTColorBox{proprieta}{o}{nota,
  before upper={\textbf{Proprietà\IfValueT{#1}{ (#1)}.}\ }}
\NewTColorBox{assioma}{o}{nota, boxrule=1pt,
  before upper={\textbf{Assioma\IfValueT{#1}{ (#1)}.}\ }}

% --- Ambienti semplici (senza riquadro) ---
% Il titolo sta in cima al blocco, anche se il contenuto inizia con un riquadro o
% con codice; e non resta mai da solo in fondo a una pagina.
% Uso: \begin{esempio}[Titolo facoltativo] ... \end{esempio}
\newcommand{\newrunin}[2]{%
  \NewDocumentEnvironment{#1}{o}{%
    \par\Needspace{4\baselineskip}\addvspace{\medskipamount}\noindent
    \textbf{#2}\IfValueT{##1}{ (##1)}\textbf{.}\ \ignorespaces}%
  {\par\addvspace{\medskipamount}}}
\newrunin{esempio}{Esempio}
\newrunin{esercizio}{Esercizio}
\newrunin{osservazione}{Osservazione}

% V e F nelle tavole di verità
\newcommand{\V}{\textsf{V}}
\newcommand{\F}{\textsf{F}}

% --- Codice C ---
% Uso: \begin{codice} ... \end{codice}      codice C numerato
%      \begin{comando} ... \end{comando}    comandi da terminale
%      \begin{risultato} ... \end{risultato} output di un programma
%      \lstinline|printf("ciao");|           codice dentro al testo
\lstdefinestyle{c}{
  language=C,
  basicstyle=\ttfamily\small,
  keywordstyle=\color{blue}\bfseries,
  commentstyle=\color{gray}\itshape,
  stringstyle=\color{red!70!black},
  numbers=left,
  numberstyle=\tiny\color{gray},
  numbersep=8pt,
  columns=flexible,
  keepspaces=true,
  breaklines=true,
  showstringspaces=false,
  tabsize=4,
  morekeywords={include, define},
  % lettere accentate nei commenti
  literate={à}{{\`a}}1 {è}{{\`e}}1 {é}{{\'e}}1 {ì}{{\`i}}1
           {ò}{{\`o}}1 {ù}{{\`u}}1 {È}{{\`E}}1 {À}{{\`A}}1
}
\lstset{style=c}
\newtcblisting{codice}{listing only, nota, left=6mm,
  listing options={style=c}}
\newtcblisting{comando}{listing only, nota,
  listing options={basicstyle=\ttfamily\small, numbers=none, language={},
    columns=flexible, keepspaces=true, breaklines=true}}
\newtcblisting{risultato}{listing only, enhanced jigsaw, breakable,
  colback=black!5, colframe=black, boxrule=0.5pt, sharp corners,
  listing options={basicstyle=\ttfamily\small, numbers=none, language={},
    columns=flexible, keepspaces=true, breaklines=true}}

% --- Diagrammi di flusso (simboli standard) ---
\tikzset{
  terminale/.style = {ellipse, draw, minimum width=2.5cm, minimum height=0.9cm},
  io/.style        = {trapezium, trapezium left angle=70, trapezium right angle=110,
                      draw, minimum width=2.5cm, minimum height=0.9cm},
  processo/.style  = {rectangle, draw, minimum width=2.5cm, minimum height=0.9cm},
  decisione/.style = {diamond, draw, aspect=2, minimum width=2.5cm},
  freccia/.style   = {thick, -{Stealth}}
}

\title{Funzioni}
\author{Marco Cerbella}
\date{Lezione 3 -- 28 settembre 2026}

\begin{document}
\maketitle

\section{Definizione di funzione}

\begin{definizione}
Siano $A$ e $B$ due insiemi. Una \emph{funzione} $f$ da $A$ a $B$ è un
sottoinsieme del prodotto cartesiano $A \times B$ tale che
\[
\forall\, a \in A \quad \exists!\, b \in B : (a, b) \in f.
\]
\end{definizione}

In altre parole, ogni elemento di $A$ compare come prima componente in
\textbf{una e una sola} coppia di $f$. Si usano le notazioni seguenti:
\begin{itemize}
    \item $f : A \to B$ indica che $f$ è una funzione da $A$ a $B$;
    \item $a \mathrel{f} b$ oppure $f(a) = b$ significa $(a, b) \in f$;
          $b$ si dice \emph{immagine} di $a$;
    \item $A$ è il \emph{dominio} di $f$ e $B$ è il \emph{codominio} di $f$.
\end{itemize}

Il sottoinsieme di $A \times B$
\[
f = \{\, (a, f(a)) \mid a \in A,\ f(a) \in B \,\} \subseteq A \times B
\]
si chiama \emph{grafico} della funzione $f$.

\begin{osservazione}
Secondo la definizione, la funzione \emph{coincide} con il suo grafico.
Le due condizioni che la definiscono sono:
\begin{enumerate}
    \item \textbf{esistenza}: ogni $a \in A$ ha almeno un'immagine in $B$
          (nessun elemento del dominio resta ``senza freccia'');
    \item \textbf{unicità}: ogni $a \in A$ ha \emph{una sola} immagine
          (da ogni elemento del dominio parte una sola freccia).
\end{enumerate}
Invece \emph{non} si chiede che ogni elemento di $B$ sia immagine di
qualcuno, né che sia immagine di uno solo: nell'esempio~2 qui sotto $a$ e
$b$ hanno ciascuno due elementi che arrivano su di loro.
\end{osservazione}

\section{Esempi}

\begin{esempio}
\leavevmode
\begin{enumerate}
    \item Siano $A = \{x, y, z\}$ e $B = \{1, 2, 3, 4\}$ e sia
          $f : A \to B$ data da
          \[
          f(x) = 1, \qquad f(y) = 3, \qquad f(z) = 4.
          \]
          Quindi $x \mathrel{f} 1$, cioè $(x, 1) \in f$, e il grafico è
          $f = \{(x,1), (y,3), (z,4)\}$. L'elemento $2 \in B$ non è
          immagine di nessun elemento di $A$.

    \item Siano $B = \{1, 2, 3, 4\}$ e $C = \{a, b\}$ e sia
          $g : B \to C$ data da
          \[
          g(1) = b, \quad g(2) = b, \quad g(3) = a, \quad g(4) = a.
          \]
\end{enumerate}

\begin{center}
\begin{tikzpicture}[scale=0.9, every node/.style={font=\small}]
  % funzione f : A -> B
  \draw (0,0) ellipse (0.9 and 1.4);
  \draw (3,0) ellipse (0.9 and 1.7);
  \node at (0,1.75) {$A$};
  \node at (3,2.05) {$B$};
  \node at (1.5,2.0) {$f$};
  \foreach \n/\y in {x/0.8, y/0, z/-0.8} {
    \fill (-0.3,\y) circle (1.5pt) node[left] {$\n$};
  }
  \foreach \n/\y in {1/1.1, 2/0.35, 3/-0.4, 4/-1.15} {
    \fill (2.7,\y) circle (1.5pt) node[right] {$\n$};
  }
  \draw[->] (-0.3,0.8) -- (2.7,1.1);
  \draw[->] (-0.3,0) -- (2.7,-0.4);
  \draw[->] (-0.3,-0.8) -- (2.7,-1.15);

  % funzione g : B -> C
  \begin{scope}[xshift=8cm]
    \draw (0,0) ellipse (0.9 and 1.7);
    \draw (3,0) ellipse (0.9 and 1.0);
    \node at (0,2.05) {$B$};
    \node at (3,1.35) {$C$};
    \node at (1.5,2.0) {$g$};
    \foreach \n/\y in {1/1.1, 2/0.35, 3/-0.4, 4/-1.15} {
      \fill (-0.3,\y) circle (1.5pt) node[left] {$\n$};
    }
    \fill (2.7,0.5) circle (1.5pt) node[right] {$b$};
    \fill (2.7,-0.5) circle (1.5pt) node[right] {$a$};
    \draw[->] (-0.3,1.1) -- (2.7,0.5);
    \draw[->] (-0.3,0.35) -- (2.7,0.5);
    \draw[->] (-0.3,-0.4) -- (2.7,-0.5);
    \draw[->] (-0.3,-1.15) -- (2.7,-0.5);
  \end{scope}
\end{tikzpicture}
\end{center}

\begin{enumerate}
    \setcounter{enumi}{2}
    \item $f : \mathbb{N} \to \mathbb{N}$, \quad $f(n) = 2n$.
          Si scrive anche $n \mapsto f(n) = 2n$.

    \item $f : \mathbb{Z} \to \mathbb{N}$, \quad
          $k \mapsto |k| =
          \begin{cases}
            k & \text{se } k \geq 0, \\
            -k & \text{se } k < 0.
          \end{cases}$

    \item $h : \mathbb{R} \to \mathbb{R}$, \quad $h(x) = x^2 - 2$.

    \item $L : \mathbb{R} \to \mathbb{R}$, \quad $L(x) = 5x + 3$.

    \item $F : \mathbb{Z} \to \mathbb{Z}$, \quad $F(k) = 3$.

    \item $F : \mathbb{Z} \to \{3\}$, \quad $F(k) = 3$.
\end{enumerate}
\end{esempio}

Gli ultimi due esempi sono funzioni \emph{costanti}: tutti gli elementi del
dominio hanno la stessa immagine. Notiamo che la legge $F(k) = 3$ è la
stessa, ma le due funzioni hanno codomini diversi ($\mathbb{Z}$ nel primo
caso, $\{3\}$ nel secondo).

\begin{definizione}
Una funzione $f : A \to B$ tale che
\[
f(x_1) = f(x_2) \qquad \forall\, x_1, x_2 \in A
\]
si dice \emph{funzione costante}.
\end{definizione}

\end{document}
